
% AirXiv conceptual primer
% Generated 2026-02-09
\documentclass[11pt]{article}

\usepackage[margin=1in]{geometry}
\usepackage{amsmath,amssymb,amsthm}
\usepackage{mathtools}
\usepackage{hyperref}
\usepackage{enumitem}

\newcommand{\Fin}{\mathrm{Fin}}
\newcommand{\NN}{\mathbb{N}}
\newcommand{\RR}{\mathbb{R}}
\newcommand{\EE}{\mathbb{E}}
\newcommand{\TV}{\mathrm{TV}}

\title{A Conceptual Primer on the Markov de Finetti Hard Direction:\\
Excursion Patterns, BEST Symmetry, and a Bridge to Universal Prediction\\
{\large (notes for formalization; 2026-02-09)}}
\author{}
\date{}

\begin{document}
\maketitle

\begin{abstract}
This note summarizes a proof architecture for the \emph{hard direction} of the Markov de Finetti
theorem (Diaconis--Freedman style): under Markov exchangeability (and an explicit recurrence
assumption), a prefix measure on sequences over a finite alphabet admits a representation as a
mixture of time-homogeneous Markov chains.

The emphasis is conceptual: we describe (i) an \emph{evidence-basis} reduction that replaces
per-word constraints by finitely many constraints on Markov transition counts; (ii) a compactness /
moment-polytope route to construct the mixing measure; and (iii) the remaining combinatorial core,
organized around \emph{excursion-pattern decompositions} and symmetry induced by \emph{segment swaps}.
The BEST theorem (Euler-trail counting) supplies the decisive uniformity mechanism.

The document is intended as a primer for future work sessions and tool-assisted proof development.
\paragraph{Status note.} Implementation-status and count claims in this draft are snapshot-scoped; verify against current repository trackers and build outputs before citing as current state. Lean-proven claims are only those tied to explicit Lean file/theorem references; broader mathematical or historical claims are literature-backed context and not asserted as Lean-proven unless explicitly stated.
\end{abstract}

\section{Executive sketch}

The overall structure can be read as a three-layer pipeline:

\begin{enumerate}[label=\textbf{L\arabic*}.]
\item \textbf{Analytic layer (parameter space + continuity).}
Represent Markov chains on a finite state space $\Fin k$ via a compact parameter space
$\Theta_k$ of initial distributions and row-stochastic transition kernels.
Cylinder probabilities $\theta \mapsto \Pr_\theta(x_{0:n})$ are continuous functions on $\Theta_k$.

\item \textbf{Evidence layer (reduce words to states).}
For each horizon $n$, the relevant statistic for Markov exchangeability is the
\emph{Markov state} (start state + transition count matrix + last state).
Let $w_\mu(n,e)$ be the prefix-measure mass of the evidence class $e$,
and $W(n,e,\theta)$ the corresponding evidence polynomial (sum of cylinder probabilities over the
fiber of trajectories with evidence $e$). The hard direction splits into:
\begin{itemize}
\item \emph{Representability:} build a probability measure $\pi$ on $\Theta_k$ such that
$w_\mu(n,e)=\int W(n,e,\theta)\,d\pi(\theta)$ for all $(n,e)$.
\item \emph{Regrouping:} prove that evidence constraints imply the per-word mixture formula
$\mu(x_{0:n})=\int \Pr_\theta(x_{0:n})\,d\pi(\theta)$.
\end{itemize}

\item \textbf{Combinatorial core (Diaconis--Freedman / BEST).}
The representability step is reduced, by compactness, to finite satisfiability of finitely many
evidence constraints. The deepest remaining ingredient is a finite-sample quantitative estimate:
an empirical (Laplace-smoothed) parameterization, averaged over long trajectories, approximates
prefix coefficients associated to short-horizon evidence. This estimate can be expressed as an
$\ell^1$ (total variation) bound between two \emph{pattern distributions}:
a with-replacement (WR) short-horizon model and a without-replacement (WOR) long-prefix model.
\end{enumerate}

The key conceptual discovery underpinning the final layer is that one should not attempt to pin a
\emph{unique} excursion multiset to a Markov state: it is generally not determined by the state.
Therefore, the final step must be phrased and proven at the \emph{distribution / fiber-sum} level,
using symmetry and averaging, rather than by selecting a canonical excursion representation.

\section{Objects and viewpoints}

\subsection{Prefix measures and Markov exchangeability}

Fix a finite alphabet (or state space) $\Fin k$.
A \emph{prefix measure} $\mu$ assigns nonnegative mass to each finite word
$x_{0:n} \in (\Fin k)^{n+1}$, consistently across horizons (a projective family).

A measure is \emph{Markov exchangeable} if its value on a word depends only on:
\begin{itemize}
\item the start state $x_0$,
\item the transition counts $N_{ab}$ for all ordered pairs $(a,b)$,
\item the last state $x_n$,
\end{itemize}
i.e.\ only on the Markov evidence. This partitions trajectories into \emph{evidence fibers}.

\subsection{The Markov parameter space}

A time-homogeneous Markov chain on $\Fin k$ is encoded by
\[
\theta = (\mathrm{init}, \mathrm{trans}), \qquad
\mathrm{init}\in \Delta(\Fin k),\quad \mathrm{trans}(a)\in \Delta(\Fin k)\ \forall a,
\]
where $\Delta(\Fin k)$ denotes probability measures on $\Fin k$.
The product topology makes $\Theta_k$ compact; cylinder probabilities
$\theta \mapsto \Pr_\theta(x_{0:n})$ are continuous.

\subsection{Evidence polynomials}

For a fixed horizon $n$ and evidence class $e$ (start, counts, last), define:
\begin{itemize}
\item $w_\mu(n,e)$: total $\mu$-mass of the evidence class $e$;
\item $W(n,e,\theta)$: the \emph{evidence polynomial} under parameter $\theta$,
the sum of $\Pr_\theta(x_{0:n})$ over all trajectories with evidence $e$.
\end{itemize}
Because $W(n,e,\theta)$ is a finite sum of continuous cylinder probabilities, it is continuous in
$\theta$.

\section{Why the hard direction can be separated into two steps}

The Markov de Finetti hard direction can be profitably factored into:

\subsection{(A) Representability at the evidence level}

Find a mixing measure $\pi$ on $\Theta_k$ such that for all $n$ and $e$,
\[
w_\mu(n,e) = \int_{\Theta_k} W(n,e,\theta)\,d\pi(\theta).
\]
This is a \emph{moment matching} problem against a basis of continuous functions $W(n,e,\cdot)$.

\subsection{(B) Regrouping from evidence constraints to word constraints}

Given such a $\pi$, show that for every word $x_{0:n}$,
\[
\mu(x_{0:n}) = \int_{\Theta_k} \Pr_\theta(x_{0:n})\,d\pi(\theta).
\]
Conceptually, this uses two constant-on-fiber facts:
\begin{itemize}
\item $\mu$ is constant on evidence fibers by Markov exchangeability;
\item for fixed $\theta$, $\Pr_\theta(\cdot)$ is constant on evidence fibers
because the Markov chain probability depends only on counts.
\end{itemize}
Then each evidence weight is a fiber cardinality times a representative probability, and the
multiplicative fiber factor cancels when equating $w_\mu$ and $\int W\,d\pi$.

Step (B) is typically routine once the evidence objects are set up; step (A) contains the real
mathematics.

\section{Compactness and the moment polytope picture}

A standard compactness argument reduces the infinite family of constraints in (A) to
\emph{finite satisfiability}:

\begin{quote}
For every finite index set $u$ of pairs $(n,e)$, there exists a probability measure $\pi$ on
$\Theta_k$ satisfying all constraints in $u$ simultaneously.
\end{quote}

Fix a finite $u$. Consider the map
\[
\pi \longmapsto \Big( \int W(n,e,\theta)\,d\pi(\theta)\Big)_{(n,e)\in u}.
\]
Its range is a compact subset of $\RR^{u}$ (or $\mathbb{R}_{\ge 0}^u$ in an extended setting),
often called a \emph{moment polytope} or \emph{moment set}. The finite satisfiability goal is
equivalent to showing the empirical constraint vector
$\big(w_\mu(n,e)\big)_{(n,e)\in u}$ lies in this compact range.

This viewpoint suggests multiple possible approaches:
\begin{itemize}
\item convex-analytic separation (show the constraint vector lies in the closed convex hull of
evaluation vectors);
\item functional-analytic extension (positive linear functionals + Riesz--Markov--Kakutani);
\item probabilistic constructions (martingales / tail sigma-fields yielding a random transition
matrix).
\end{itemize}
In the Diaconis--Freedman route, the recurrence hypothesis enables explicit empirical
approximations that realize the constraint vector as a limit of points in the moment polytope.

\section{Empirical measures, recurrence, and smoothing}

A central construct is an \emph{empirical parameter map} from evidence states to Markov parameters.
Given a Markov state $s$ (with transition counts), define a Laplace-smoothed empirical transition
kernel
\[
\widehat{P}_s(b\mid a)=\frac{N_{ab}(s)+\alpha_b}{N_{a\bullet}(s)+\sum_{b'}\alpha_{b'}},
\]
with a symmetric prior (e.g.\ $\alpha_b=1$). This defines $\theta_s$.

A long-horizon empirical mixing measure is then obtained by pushing forward the distribution of
evidence states at horizon $N$ through $s\mapsto \theta_s$.
The main convergence goal is:

\begin{quote}
For each fixed short horizon $n$ and evidence class $e$, the empirical integral
$\int W(n,e,\theta)\,d\pi_N(\theta)$ converges to $w_\mu(n,e)$ as $N\to\infty$.
\end{quote}

Recurrence (returns to a start/anchor state) supplies enough regeneration to control the empirical
approximation error. The remaining error is quantified by a bound of order
$O((n+1)^2/R)$, where $R$ is the number of returns to the anchor state in the long-horizon evidence.

\section{Excursion patterns: turning the core bound into a finite combinatorial problem}

\subsection{Excursions and patterns}

Fix an anchor state (often the start state).
A long trajectory can be cut at return times to the anchor into \emph{excursions}.
For many arguments, the key object is the \emph{excursion list} (a list of excursion segments)
and its derived \emph{pattern} information.

Two distributions over patterns are compared:

\begin{itemize}
\item \textbf{WR (with replacement):} sample short-horizon trajectories (length $n+1$) and weight
each by its Markov chain probability under the empirical parameter $\theta_s$.

\item \textbf{WOR (without replacement):} sample a long-horizon trajectory consistent with evidence
$s$ uniformly within its evidence fiber, take its length-$n+1$ prefix, and map to an excursion
pattern.
\end{itemize}

The difference is measured by an $\ell^1$ sum over a finite pattern index set:
\[
\sum_{p} \big| \mathrm{WR}(p)-\mathrm{WOR}(p)\big|.
\]
The goal is to bound this sum by $4(n+1)^2/R$ (up to constants).

\subsection{A crucial modeling caution}

A tempting shortcut is to assume the evidence state determines a unique multiset of excursions.
In general, it does not. Evidence (transition counts + endpoints) constrains the multigraph of
transitions, but multiple distinct collections of excursion segments may realize the same counts.

Therefore, any proof step that ``fixes the excursion multiset'' must be justified as a conditioning
argument (i.e.\ prove a uniform bound that is stable under averaging), or avoided.
The robust alternative is to work directly with finite sums over patterns and to use symmetry
mechanisms (segment swaps, adjacent transpositions) to show certain terms are constant on
orbits/fibers.

\section{Segment swaps and symmetry transfer}

\subsection{Local symmetry operation}

A \emph{segment swap} is an operation on a trajectory that exchanges two adjacent segments of
lengths $L_1,L_2$ starting at position $a$, producing a new trajectory:
\[
x = \mathrm{pre} \cdot u \cdot v \cdot \mathrm{suf}
\quad\mapsto\quad
x'=\mathrm{pre} \cdot v \cdot u \cdot \mathrm{suf}.
\]
Under appropriate boundary conditions (notably, returns to the anchor at the cut points),
a segment swap preserves the overall evidence state while permuting excursion order.

\subsection{Fiber-image equalities}

A recurring pattern in the formalization is:
\begin{quote}
If segment swap maps one pattern fiber into another and vice versa, then the two fibers have
equal cardinality, hence equal ratio masses.
\end{quote}
This is packaged as a reusable lemma schema:
\begin{itemize}
\item image equality $\Rightarrow$ ratio invariance for WOR-side mass;
\item image equality $\Rightarrow$ ratio invariance for short-horizon ratios;
\item combined invariance $\Rightarrow$ invariance of the absolute discrepancy term.
\end{itemize}

Conceptually, this converts a local bijection into equality of contributions in the target $\ell^1$
sum.

\section{Orbit / fiber decompositions via multisets}

\subsection{Why partition by multiset}

Many invariances act only within the set of patterns that share the same multiset of excursions.
A natural way to organize the global sum is therefore to partition the pattern index set by
\[
p \longmapsto \mathrm{Multiset}(p),
\]
i.e.\ map an excursion list to the multiset of its elements (forgetting order).

This converts the global discrepancy sum into a sum of inner sums over ``multiset fibers'':
\[
\sum_{\mathcal{M}} \ \sum_{p:\,\mathrm{Multiset}(p)=\mathcal{M}} \big|\mathrm{WR}(p)-\mathrm{WOR}(p)\big|.
\]
Inside a fixed fiber, adjacent transpositions generate all permutations of the excursion list, so
one can aim to prove that the summand is \emph{constant} on the fiber (or at least controlled by a
simple expression).

\subsection{From adjacent swaps to arbitrary permutations}

A common proof pattern is:

\begin{enumerate}
\item show invariance under a single adjacent transposition (via a segment swap lemma);
\item lift to invariance under any permutation using a permutation induction
(e.g.\ decomposing permutations into adjacent swaps);
\item conclude that the function is constant on each multiset fiber.
\end{enumerate}

The technical work is (1): producing the right decomposition witnesses and boundary conditions so
that the segment swap lemma applies to the chosen adjacent transposition. Once (1) is available,
(2)--(3) are largely algebraic.

\section{Where the BEST theorem enters}

The BEST theorem is a counting formula for Euler trails in a directed multigraph.
In this setting, a Markov evidence state specifies a directed multigraph whose edge multiplicities
are the transition counts. Trajectories in an evidence fiber correspond to Euler trails (with a
vertex-sequence extraction map).

The conceptual connection is:

\begin{quote}
Counting Euler trails supplies an explicit formula for evidence-fiber cardinalities and for
sub-fiber cardinalities under ``fixing part of the excursion order''. This yields quantitative
control of how the prefix distribution deviates from the idealized with-replacement model.
\end{quote}

Informally, if the long trajectory has $R$ returns to the anchor, then the first $m=n+1$ segments
behave almost as if sampled with replacement from a large pool of $R$ excursions; the deviation is
of order $m^2/R$. The BEST theorem makes the ``pool symmetry'' exact enough to turn this intuition
into a bound.

\section{A practical proof template (tool-friendly)}

This section outlines a proof-template that is robust under refactoring and suitable for formal
proof assistants.

\subsection{Isolation of the core lemma}

Treat the final quantitative bound as a single named lemma (the ``core'') and make everything else
feed into it. The core lemma should:
\begin{itemize}
\item quantify over $k,n,N,e,s$ and include the recurrence/return-size hypotheses;
\item rewrite the target discrepancy into a finite sum over a canonical pattern index set;
\item apply a partition-by-multiset identity so the remaining goal is purely fiberwise.
\end{itemize}

\subsection{Symmetry layer}

Prove reusable symmetry lemmas in the following order:
\begin{enumerate}
\item \textbf{Local swap:} adjacent transposition on excursion order realized as a segment swap
on trajectories (requires bridge lemmas about return positions).
\item \textbf{Fiber mapping:} show segment swap maps the appropriate pattern fibers to each other.
\item \textbf{Term invariance:} package invariance of the discrepancy summand.
\item \textbf{Permutation lift:} extend invariance from adjacent swaps to arbitrary permutations
inside a multiset fiber.
\end{enumerate}

\subsection{Counting / bounding layer}

On each multiset fiber, use Euler-trail counting (BEST) to compare:
\begin{itemize}
\item the distribution induced by uniform choice in a long evidence fiber (WOR),
\item the distribution induced by Markov probabilities on short evidence (WR),
\end{itemize}
and establish the $O(m^2/R)$ bound, then sum over fibers.

A pragmatic note: constants can often be loosened early (prove $C\cdot m^2/R$ with a generous
$C$) and tightened later if needed.

\section{Generalizations and ``better indexing''}

Once the proof is established, a useful conceptual cleanup is to re-index the theorem around an
\emph{anchor state} rather than a random initial state. In many Markov formalizations, one can
assume an initial state is fixed almost surely; when not, one can often condition on an anchor and
lift results back to the mixture statement. This re-indexing improves composability with other
systems (e.g.\ prediction frameworks) because it makes the regeneration structure explicit.

\section{Summary}

The Markov de Finetti hard direction is best understood as:
\begin{enumerate}
\item a compact parameter space + continuous cylinder kernels;
\item an evidence basis that reduces per-word constraints to moment constraints;
\item a compactness reduction to finite satisfiability;
\item an empirical approximation route enabled by recurrence;
\item a final finite combinatorial inequality, naturally expressed as a discrepancy between WR and
WOR \emph{excursion-pattern distributions}, discharged via segment-swap symmetries and BEST/Euler-trail
counting.
\end{enumerate}

The most durable methodological lesson is to keep the last combinatorial inequality isolated and
prove it with orbit/fiber symmetry, avoiding assumptions (such as a uniquely determined excursion
multiset) that are not valid at the evidence-state level.

\section*{Pointers for future work}

\begin{itemize}
\item Prove (or reuse) a generic ``with- vs without-replacement'' total variation lemma with an
$O(m^2/R)$ bound, and instantiate it to the excursion-pattern distributions. This can reduce the
dependency on bespoke orbit bookkeeping.
\item Where the formalization uses Laplace smoothing, track the smoothing error explicitly:
either show it cancels in the ratio form, or add it as an additive $O(m/R)$ correction.
\item Consider modularizing the BEST/Euler-trail connection as its own component: a library lemma
that turns Euler-trail counting into near-uniformity of prefix distributions under recurrence.
\end{itemize}

\bibliographystyle{plain}
\begin{thebibliography}{9}
\bibitem{DF80}
P.~Diaconis and D.~Freedman.
\newblock De Finetti's theorem for Markov chains.
\newblock \emph{Annals of Probability}, 8(1):115--130, 1980.

\bibitem{BEST}
N.~de Bruijn, T.~van Aardenne-Ehrenfest, C.~A. Smith, and W.~Tutte.
\newblock The BEST theorem (Euler trail counting) and related combinatorics.
\newblock Standard graph theory references; see also modern expositions.

\bibitem{Solomonoff}
R.~J. Solomonoff.
\newblock A formal theory of inductive inference.
\newblock \emph{Information and Control}, 1964.

\end{thebibliography}

\end{document}
